\documentclass[11pt,a4paper]{article}
\usepackage[T1]{fontenc}
\usepackage{lmodern,microtype}
\usepackage[hmargin=25mm,vmargin=23mm]{geometry}
\usepackage{amsmath,amssymb,booktabs}
\usepackage{xurl}
\usepackage[hidelinks]{hyperref}
\hypersetup{
  pdftitle={ESDM v2.0.1: Graph-informed encoder initialisation},
  pdfauthor={Lorenzo Moriondo}}
\title{\textbf{ESDM v2.0.1: Graph-informed encoder initialisation}\\
  \large A synthetic follow-up to v2.0.0}
\author{Lorenzo Moriondo\\[-1pt]
  \small tuned.org.uk \quad\textbar\quad
  \href{https://orcid.org/0000-0002-8804-2963}{ORCID 0000-0002-8804-2963}}
\date{21 November 2026}

\begin{document}
\maketitle
\begin{abstract}
Initialising the posterior-mean encoder from an existing graph chart
improves clean-feature affine recoverability on synthetic smooth path
signals. A frozen 25-candidate search selected 16 updates at learning rate
$3\times10^{-5}$. On 2048 fresh confirmation cases, mean squared error
fell from $0.09331875$ for the default-budget random control to
$0.04582097$, a 50.898\% reduction. The reduction against a budget-matched
random control was 48.449\%. However, the selected model was
indistinguishable from its untrained chart-initialised counterpart under
the paired case interval. The evidence supports an informative
initialisation, not improved representation learning through training.
This addendum documents the v2.0.1 follow-up; the original v2.0.0 manuscript
and its historical results remain available at that release tag.
\end{abstract}

\section{The intervention}
The Entropic Semantic Diffusion Model (ESDM) already retains an orthonormal
chart $U_q$ consisting of the lowest-frequency eigenvectors of its frozen
feature Laplacian. Version 2.0.1 uses that chart to initialise the existing
trainable posterior-mean layer. With row-vector observations,
\[
  \mu(x)=xW_\mu^\top+b_\mu,\qquad
  W_\mu^{(0)}=U_q^\top,\qquad b_\mu^{(0)}=0.
\]
In \texttt{esdm/models/esdm.py}, this is a copy into
\texttt{enc\_mu.weight} and a zeroing of \texttt{enc\_mu.bias}, performed
without gradient tracking at construction. The layer remains trainable;
it is not tied to the frozen chart thereafter. No parameters, dimensions,
loss terms or inference constraints are added. The log-variance head and
Gaussian sampling rule retain their previous initialisation and behaviour.
Latent samples and vibrational amplitudes remain free, while reconstruction
modal gains retain the v2.0.0 bounds. Checkpoint loading replaces the
initial weights normally.

\section{Frozen search and measurement}
The experiment follows \texttt{autoresearch.md}, with v2.0.0 commit
\texttt{c674a5c} as the baseline. It uses an unweighted 32-node path and
rank $q=8$. Clean signals are drawn in its analytic discrete cosine
transform (DCT-II) eigenbasis $U$:
\[
  x=aU^\top,\qquad
  a_k\sim\mathcal N\!\left(0,\frac{1}{1+(k/4)^2}\right),
  \quad k=0,\ldots,31,
\]
with independent coefficients. Consequently the graph prior is informative
by construction. Both model training and probe evaluation use clean
signals; the observation noise used in the earlier readout experiment is
not used here.

All runs use CPU execution and seed 3407. NumPy
\texttt{SeedSequence} spawn keys distinguish 512 training cases (stream 0),
512 fresh validation cases (100), and 2048 fresh confirmation cases (101).
The old stream-1 test set is development evidence only. The search crosses
five learning rates, $3\times10^{-5}$, $10^{-4}$, $3\times10^{-4}$,
$10^{-3}$ and $3\times10^{-3}$, with maximum update counts
16, 32, 64, 128 and 256. These are experimental budgets, not new model
constants. Every candidate trains the full model using its actual ESDM
loss, Adam, weight decay $10^{-5}$, batch size 64 and gradient clipping
at 1.0. Stream 102 supplies shared epoch permutations; model and training
randomness are reset identically across runs. Training stops at heat death
or the budget; all 25 recorded candidates reached their budgets.

The primary metric is mean squared error (MSE) per original feature from
an affine least-squares probe of the posterior mean. Each probe is fitted
only on the 512 training cases, including an intercept, and then evaluated
without refitting. This measures linear recoverability, not the model's
forward reconstruction, diffusion quality or generation. The selected
candidate minimises validation MSE; exact ties favour fewer updates and
then a lower rate. Its checkpoint is saved before confirmation cases are
generated, with no subsequent retuning.

Two additional training runs provide random-mean controls: the baseline
rate $3\times10^{-4}$ with 32 updates, and the selected rate and budget.
Their mean weights and biases are restored from the frozen v2.0.0
initialisation; other initial parameters are shared. The baseline reported
below is this protocol's control, not a relabelling of the original
manuscript's differently batched diagnostic.

Acceptance requires at least a 1\% confirmation MSE reduction against
\emph{each} trained random control, with the upper endpoint of each paired
95\% interval below zero. Differences are selected-model minus control
per-case errors. The intervals are percentile case-bootstrap intervals
from 2000 resamples of the 2048 paired cases (stream 103), without probe
refitting. They are conditional on one training seed, the fitted models
and probes, and the fixed training data. They do not quantify training-seed
uncertainty and are not simultaneous, multiplicity-adjusted intervals.

\section{Confirmation results}
The selected candidate used $3\times10^{-5}$ and 16 updates, obtaining
validation MSE $0.04509079082$. Table~\ref{tab:mse} reports confirmation
errors; Table~\ref{tab:contrasts} gives the paired contrasts. Both random
controls satisfy the frozen acceptance rule, with relative reductions of
50.898426\% and 48.448689\%, respectively.

\begin{table}[ht]
\centering
\small
\begin{tabular}{@{}lrr@{}}
\toprule
Representation or reconstruction & Updates & Confirmation MSE\\
\midrule
Random mean, baseline rate & 32 & 0.09331874676\\
Random mean, matched rate & 16 & 0.08888420540\\
Chart-initialised mean, selected & 16 & 0.04582097311\\
Chart-initialised mean, untrained & 0 & 0.04582100595\\
Fixed chart, transpose synthesis & 0 & 0.04505789528\\
\bottomrule
\end{tabular}
\caption{The first four rows use separately fitted training-only affine
probes. The last reconstructs $xU_qU_q^\top$ directly, with no fitted
probe, and is a distinct estimator.}
\label{tab:mse}
\end{table}

\begin{table}[ht]
\centering
\small
\begin{tabular}{@{}lrr@{}}
\toprule
Comparator & Mean difference & Paired 95\% interval\\
\midrule
Baseline random & $-0.04749777365$ & $[-0.04929368082,\,-0.04575577022]$\\
Matched random & $-0.04306323229$ & $[-0.04475371834,\,-0.04142955549]$\\
Untrained chart & $-3.2835133\times10^{-8}$ &
$[-2.3961762,\,1.6180500]\times10^{-7}$\\
\bottomrule
\end{tabular}
\caption{Negative differences favour the selected model. The untrained-chart
contrast is an attribution check, not an acceptance comparator.}
\label{tab:contrasts}
\end{table}

The selected and untrained chart means differ by only
$-3.28\times10^{-8}$ MSE, with an interval crossing zero. This does not
prove exact equivalence, but supplies no evidence that training improved
on the initial chart. Direct fixed-chart synthesis has still lower error,
$0.04505789528$, without a fitted decoder. It should not be conflated with
the train-only affine-probe estimator.

Every other chart candidate had worse validation error than the selected
one. At 16 updates, increasing the rate to $0.003$ raised MSE to
$0.04515054$; at the selected rate, extending to 256 updates raised it to
$0.04509290$. More aggressive training could be much worse: rate $0.003$
gave $0.12535620$ at 128 updates and $0.11977122$ at 256. Thus longer
training generally degraded recovery here, but the dependence was not
strictly monotonic.

\section{Interpretation and audit trail}
Affine probes can absorb invertible latent rescaling. The selected
encoder remained rank eight, with singular values between 0.997195 and
0.998187, versus 0.339419--0.754769 for the baseline random control.
Selected posterior-mean variances ranged from 0.233299 to 1.003858 per
mode; baseline values ranged from 0.045068 to 0.073672. The corresponding
posterior Kullback--Leibler terms were 2.552077 and 0.365147. These
diagnostics accompany the probe result rather than turning it into a
claim about calibrated latent sampling.

A separately fitted affine probe of one posterior draw per case obtained
MSE 0.139594 for the selected model, 0.194469 for the baseline and 0.192391
for the matched control. The untrained chart obtained 0.139448. These are
descriptive sampled-code errors; no confidence intervals for these
comparisons were recorded. They do not establish a generation benefit.
The selected training trace records minibatch loss 19.068607 at the first
update and 18.146175 at the last; different minibatches prevent treating
these endpoints as a controlled before--after loss comparison.

The defensible interpretation is retention of a useful low-frequency
prior. Changing initialisation produces a large recoverability gain over
random means on this deliberately smooth distribution, while optimisation
adds no demonstrated gain over the untrained chart. Neither broader
semantic representation quality, an advantage of the prime-based gain
initialisation, nor image or audio generation is tested by this follow-up.

The frozen protocol, selection, all 25 candidate scores and confirmation
statistics are stored under
\path{autoresearch/loop-261121-representation/} as
\texttt{protocol.json}, \texttt{selection.json}, \texttt{results.tsv} and
\texttt{confirmation.json}; \texttt{autoresearch.jsonl} retains the
training traces and both controls. The measured source snapshot is
\texttt{55e830d17af8}, with source and data hashes recorded in the protocol;
the selected checkpoint hash begins \texttt{bb55fe129d8b}. The measured
runtime was PyTorch 2.12.1 and NumPy 2.5.0, distinct from the updated release
environment. A subsequent frozen-checkpoint compatibility replay under
PyTorch 2.15.0 and NumPy 2.5.0 reproduced all three primary MSEs exactly
(difference zero), as recorded in \texttt{runtime-validation.json}.
That replay involved no retraining, reselection or new acceptance test;
the reported search measurements and intervals remain those of the original
runtime. Later TSV newline normalisation is output formatting only; the
protocol's source hashes identify the measured snapshot.

The runner checks finite losses,
gradients and states, exact preservation of the frozen chart, and bounded
readout gains. Confirmation records exact equality of state tensors and
validation errors after checkpoint restoration. These records are
distinct from the protocol's separately required full-test and lint
release checks.

The archived \path{docs/paper/esdm-v2.tex} at tag
\href{https://github.com/tuned-org-uk/entropic-semantic-diffusion/tree/v2.0.0}{v2.0.0}
remains the original manuscript. This separate v2.0.1 addendum reports a
synthetic initialisation follow-up, without replacing its historical
results.
\end{document}
